\documentclass[11pt]{article}

\usepackage[margin=1in]{geometry}
\usepackage{amsmath,amssymb,amsthm,mathtools}
\usepackage{graphicx}
\usepackage{microtype}
\usepackage[hidelinks]{hyperref}
\usepackage{xcolor}
\setlength{\emergencystretch}{2em}

\newtheorem{theorem}{Theorem}
\newtheorem{lemma}[theorem]{Lemma}
\newtheorem{proposition}[theorem]{Proposition}
\newcommand{\R}{\mathbb R}
\newcommand{\Hh}{\mathbb H}
\newcommand{\B}{\mathcal B}
\newcommand{\A}{\mathrm A}
\newcommand{\VN}{\mathrm{VN}}
\newcommand{\id}{\mathrm{id}}
\newcommand{\dd}{\,d}
\newcommand{\wt}{\widetilde}
\newcommand{\PhiA}{\Phi_{\A(\pi_+)}}

\title{The Coefficient-Algebra Spectrum of the Real Affine Group\\
\large A full candidate solution to the $ax+b$ conjecture in arXiv:1112.4878}
\author{Automated mathematics run \texttt{fa\_banach\_001}, lane 1}
\date{26 August 2026}

\begin{document}
\maketitle

\begin{center}
\fbox{\parbox{0.91\textwidth}{\textbf{Status.}
Full solution candidate, likely valid, awaiting human review.  The packet
proves the concrete conjecture in Spronk--Stokke's $ax+b$ example: the
coefficient-norm algebra generated by the positive irreducible representation
has no characters beyond the closed upper-half-plane complexification found
for its uniform closure.  This proves spectral naturality for this
representation; it does not prove the paper's broader Conjecture~3.9(ii).}}
\end{center}

\section{Source problem}

The source is N.~Spronk and R.~Stokke, \emph{Matrix coefficients of unitary
representations and associated compactifications}, arXiv:1112.4878, published
in \emph{Indiana University Mathematics Journal} \textbf{62} (2013),
99--148 \cite{SpronkStokke}.  Its final example studies the real affine group
\[
 G=\{(a,b):a>0,\ b\in\R\},\qquad
 (a,b)(a',b')=(aa',ab'+b),
\]
and the irreducible representation on
$H=L^2(\R_{>0},\dd s)$
\begin{equation}\label{eq:pi}
 (\pi_+(a,b)f)(s)=a^{1/2}e^{ibs}f(as).
\end{equation}
Write $\A(\pi_+)$ for the closure, in Fourier--Stieltjes norm, of the
algebra generated by the matrix coefficients of $\pi_+$.  The source defines
\[
 \wt G=\{(a,z):a>0,\ z\in\Hh\},\qquad
 \Hh=\{z\in\mathbb C:\operatorname{Im}z\geq0\},
\]
with $(a,z)(a',z')=(aa',az'+z)$, and extends \eqref{eq:pi} to contractions
\begin{equation}\label{eq:tpi}
 (\wt\pi_+(a,z)f)(s)=a^{1/2}e^{izs}f(as).
\end{equation}
On page 37, after checking that these contractions are characters, it asks
whether they are all the characters.

\begin{figure}[ht]
\centering
\includegraphics[width=0.99\textwidth]{figures/open_problem_crop.png}
\caption{The operator criterion, the complexified semigroup, and the conjecture
in the source paper, PDF page 37.}
\end{figure}

The exact conjecture is
\[
 \boxed{\ \PhiA\cong\wt G.\ }
\]
The source also proves that the spectrum of the corresponding uniform algebra
is $\wt G$.  Thus its conjecture says exactly that $\pi_+$ is spectrally
natural.

\section{Proof intuition}

The tensor square of $\pi_+$ becomes very transparent in sum--ratio
coordinates
\[
 r=s+t,\qquad u=\frac{s}{s+t}.
\]
It is just $\pi_+$ in the radial variable $r$, amplified by the identity in
the angular variable $u$.  Consequently a character operator $x$ satisfies
\[
 \Delta(x)=x\otimes x,
\]
where $\Delta$ is radial amplification.  The right side must therefore
commute with every multiplier depending only on $s/(s+t)$.

For a positive group-like operator $T$, this angular commutation is rigid.
If $T$ had an off-diagonal block between two separated radial intervals, one
could tensor it with a nonzero positive diagonal block of $T$ on a sufficiently
narrow third interval.  The resulting block of $T\otimes T$ would connect two
disjoint angular ranges, contradicting angular commutation.  Hence $T$ is a
multiplication operator.  Its group-like equation is then the measurable
Cauchy equation $f(s+t)=f(s)f(t)$, so $T=M_{e^{-cs}}$, $c\geq0$.

For a unitary group-like operator $V$, angular commutation preserves the
operator $Q_1(Q_1+Q_2)^{-1}$, where $Q=M_s$.  Applying the bounded imaginary
powers of the logit function shows that
$V^*QV=\alpha Q$.  Therefore $V$ is a dilation followed by a unimodular
multiplier.  The group-like equation again becomes the Cauchy equation and
forces that multiplier to be $e^{ibs}$.  Thus $V=\pi_+(a,b)$.

Finally, the polar decomposition of an arbitrary nonzero character is again
group-like.  The positive and unitary classifications give
\[
 x=\pi_+(a,b)M_{e^{-cs}}=\wt\pi_+(a,b+iac),
\]
which is exactly the claimed closed-upper-half-plane family.

\section{The radial coproduct}

Let $K=L^2((0,1),\dd u)$ and define a unitary
$W:H\otimes H\to L^2(\R_{>0}\times(0,1),\dd r\dd u)=H\otimes K$ by
\begin{equation}\label{eq:W}
 (W\xi)(r,u)=r^{1/2}\xi(ru,r(1-u)).
\end{equation}
The Jacobian of $(r,u)\mapsto(ru,r(1-u))$ is $r$, so $W$ is indeed unitary.
A direct calculation gives
\begin{equation}\label{eq:intertwine}
 W(\pi_+\otimes\pi_+)(a,b)W^*=\pi_+(a,b)\otimes I_K.
\end{equation}
In particular $\pi_+\otimes\pi_+$ is an amplification of $\pi_+$, so the
closed coefficient space of $\pi_+$ is already an algebra and equals
$\A(\pi_+)$.  The
normal $*$-homomorphism occurring in the Walter criterion of
Spronk--Stokke is
\begin{equation}\label{eq:delta}
 \Delta(x)=W^*(x\otimes I_K)W\qquad(x\in\B(H)).
\end{equation}
That criterion says
\begin{equation}\label{eq:grouplike}
 x\in\PhiA\cup\{0\}
 \quad\Longleftrightarrow\quad
 \Delta(x)=x\otimes x.
\end{equation}

Let
\[
 R=\frac{Q\otimes I}{Q\otimes I+I\otimes Q},
 \qquad (Qf)(s)=sf(s).
\]
Under $W$, the operator $R$ is $I_H\otimes M_u$.  Thus \eqref{eq:delta}
commutes with every bounded Borel function of $R$.  We record the consequence
used repeatedly:
\begin{equation}\label{eq:angular}
 \Delta(x)=x\otimes x
 \quad\Longrightarrow\quad
 [x\otimes x,\psi(R)]=0
 \quad(\psi\in L^\infty(0,1)).
\end{equation}
Also $\Delta$ is isometric.  Taking norms in \eqref{eq:grouplike} shows that
every nonzero group-like operator has norm one.

\section{Positive group-like operators}

\begin{lemma}[Angular block rigidity]\label{lem:positive}
Let $0\neq T\in\B(H)$ be positive.  If $T\otimes T$ commutes with every
bounded Borel function of $R$, then $T$ is a multiplication operator.
Consequently, if in addition $\Delta(T)=T\otimes T$, then
\[
 T=M_{e^{-cs}}
\]
for a unique $c\geq0$.
\end{lemma}

\begin{proof}
For a Borel set $E\subset\R_{>0}$ write $P_E=M_{1_E}$.  First take bounded
open intervals $I=(i_0,i_1)$ and $J=(j_0,j_1)$ with $i_1<j_0$.  We prove
$P_ITP_J=0$.

Choose $q$ with $1<q<j_0/i_1$ and partition $\R_{>0}$, up to endpoints, into
the intervals $K_n=(q^n,q^{n+1})$.  At least one $K=K_n$ satisfies
\begin{equation}\label{eq:diagblock}
 P_KTP_K\neq0.
\end{equation}
Indeed, if every such compression vanished, positivity would give
$T^{1/2}P_K=0$ for every $K$, hence $T=0$.

On $I\times K$, the angular coordinate $s/(s+t)$ is at most
$i_1/(i_1+q^n)$.  On $J\times K$, it is at least
$j_0/(j_0+q^{n+1})$.  The choice $q<j_0/i_1$ makes the first number strictly
smaller than the second.  Hence there is a Borel set $S\subset(0,1)$ such
that the angular projection $E=1_S(R)$ is the identity on
$L^2(I\times K)$ and zero on $L^2(J\times K)$.  Since $T\otimes T$ commutes
with $E$,
\begin{align*}
 0
 &=(P_I\otimes P_K)E(T\otimes T)(I-E)(P_J\otimes P_K)\\
 &=(P_ITP_J)\otimes(P_KTP_K).
\end{align*}
By \eqref{eq:diagblock}, $P_ITP_J=0$.  Self-adjointness gives the reverse
off-diagonal block as well.

Exhausting $(0,t)$ and $(t,\infty)$ by separated bounded intervals and taking
strong limits now shows
\[
 P_{(0,t)}TP_{(t,\infty)}=P_{(t,\infty)}TP_{(0,t)}=0
 \qquad(t>0).
\]
Thus $T$ commutes with all spectral projections of $Q$.  The multiplication
algebra on $L^2(\R_{>0})$ is maximal abelian, so $T=M_f$ for a bounded
measurable $f\geq0$.

If $T$ is group-like, \eqref{eq:W} gives
\[
 \Delta(M_f)=M_{f(s+t)},\qquad
 M_f\otimes M_f=M_{f(s)f(t)}.
\]
Therefore $f(s+t)=f(s)f(t)$ for almost every $(s,t)$.  The standard
measurable version of the exponential Cauchy equation (obtained, if needed,
by changing $f$ on a null set) gives either $f=0$ almost everywhere or
$f(s)=e^{\kappa s}$.  The first case is excluded.  Since every nonzero
group-like operator has norm one, $\kappa\leq0$.  Writing $c=-\kappa$ proves
the assertion and uniqueness.
\end{proof}

An immediate useful consequence is that the only nonzero group-like
projection is $I$.

\section{Unitary group-like operators}

\begin{lemma}[Logarithmic ratio rigidity]\label{lem:unitary}
If $V\in\B(H)$ is unitary and $\Delta(V)=V\otimes V$, then there are unique
$a>0$ and $b\in\R$ such that
\[
 V=\pi_+(a,b)=M_{e^{ibs}}D_a,
 \qquad (D_af)(s)=a^{1/2}f(as).
\]
\end{lemma}

\begin{proof}
Put $A=V^*QV$.  From \eqref{eq:angular}, conjugation by $V\otimes V$ fixes
$R$, and hence
\begin{equation}\label{eq:ratioA}
 \frac{A\otimes I}{A\otimes I+I\otimes A}
 =\frac{Q\otimes I}{Q\otimes I+I\otimes Q}.
\end{equation}
The function $h_t(r)=(r/(1-r))^{it}$ is bounded on $(0,1)$ for each
$t\in\R$.  Applying bounded Borel functional calculus to
\eqref{eq:ratioA} gives
\begin{equation}\label{eq:imaginary}
 A^{it}\otimes A^{-it}=Q^{it}\otimes Q^{-it}.
\end{equation}
After multiplying on the left by $Q^{-it}\otimes Q^{it}$, this has the form
$C_t\otimes D_t=I$.  If two nonzero bounded operators have tensor product
equal to the identity, each is scalar; for example, slice by a vector
functional nonzero on $D_t$.  Hence
\[
 Q^{-it}A^{it}=c(t)I,
 \qquad A^{it}=c(t)Q^{it},
\]
where $|c(t)|=1$.  Strong continuity and the group laws show that $c$ is a
continuous character of $\R$, so $c(t)=e^{i\gamma t}$ for some
$\gamma\in\R$.  Uniqueness of generators of one-parameter unitary groups
then yields
\[
 A=e^\gamma Q.
\]

Set $a=e^{-\gamma}$.  Since $D_a^*QD_a=Q/a=e^\gamma Q$, the unitary
$U=VD_a^*$ commutes with $Q$.  Maximal abelianness of the multiplication
algebra gives $U=M_\varphi$ for a measurable unimodular function $\varphi$.
Thus $V=M_\varphi D_a$.

Both $V$ and $D_a=\pi_+(a,0)$ are group-like.  Cancelling
$D_a\otimes D_a$ from their group-like equations gives
\[
 \varphi(s+t)=\varphi(s)\varphi(t)
 \quad\text{for almost every }(s,t).
\]
A measurable circle-valued solution of this Cauchy equation is
$\varphi(s)=e^{ibs}$ for a unique $b\in\R$.  This proves the result.
\end{proof}

\section{Completion of the character classification}

\begin{theorem}[The $ax+b$ coefficient spectrum]\label{thm:main}
For the representation \eqref{eq:pi}, every nonzero solution of the Walter
equation \eqref{eq:grouplike} is uniquely of the form
\[
 \wt\pi_+(a,z),\qquad a>0,\quad \operatorname{Im}z\geq0.
\]
Consequently
\[
 \Phi_{\A(\pi_+)}\cong
 \wt G=\{(a,z):a>0,\ \operatorname{Im}z\geq0\}
\]
as a locally compact topological semigroup.  In particular $\pi_+$ is
spectrally natural.
\end{theorem}

\begin{proof}
Let $0\neq x$ satisfy \eqref{eq:grouplike}, and write its polar decomposition
as $x=v|x|$.  Normal $*$-homomorphisms preserve absolute values, and
\[
 |x\otimes x|=|x|\otimes|x|.
\]
Therefore $|x|$ is group-like.  The polar-decomposition part of the Walter
criterion in \cite[Theorem 3.6]{SpronkStokke} also gives that $v$ is
group-like.  Lemma~\ref{lem:positive} gives
\[
 |x|=M_{e^{-cs}}\qquad(c\geq0).
\]
The initial and final projections $v^*v$ and $vv^*$ are nonzero group-like
projections.  The consequence following Lemma~\ref{lem:positive} says both
are $I$, so $v$ is unitary.  Lemma~\ref{lem:unitary} now gives
$v=\pi_+(a,b)$ for some $a>0$, $b\in\R$.  Hence
\begin{align*}
 (xf)(s)
 &=\bigl(M_{e^{ibs}}D_aM_{e^{-cs}}f\bigr)(s)\\
 &=a^{1/2}e^{ibs}e^{-cas}f(as)
 =\wt\pi_+(a,b+iac)f(s).
\end{align*}
Since $b+iac\in\Hh$, every nonzero group-like operator belongs to the stated
family.  Conversely, substituting \eqref{eq:tpi} in \eqref{eq:W} directly
gives
\[
 \Delta(\wt\pi_+(a,z))
 =\wt\pi_+(a,z)\otimes\wt\pi_+(a,z),
\]
as already observed in the source.  Uniqueness follows first from support
transport, which determines $a$, and then from equality of the exponential
multipliers, which determines $z$.

It remains only to record the topology implicit in $\cong$.  The map
\begin{equation}\label{eq:parammap}
 (a,z)\longmapsto\wt\pi_+(a,z)
\end{equation}
is weak-operator continuous.  Adjoin one point $\infty$ to $\wt G$ and send
it to the zero operator.  This extension is continuous: for compactly
supported test vectors bounded away from zero, coefficients vanish identically
when $a\to0$ or $a\to\infty$, vanish by domination when
$\operatorname{Im}z\to\infty$, and vanish by the Riemann--Lebesgue lemma when
$|\operatorname{Re}z|\to\infty$ while the other parameters stay in compact
sets.  Density and contractivity remove the support restriction.

The set of solutions to \eqref{eq:grouplike} in the operator unit ball is
weak-operator closed: $\Delta$ is weak-operator continuous on bounded sets,
and $x\mapsto x\otimes x$ is weak-operator continuous there on elementary
matrix coefficients.  It is therefore compact.  The extended map from the
one-point compactification of $\wt G$ is, by the classification above, a
continuous bijection onto this compact solution set.  It is a homeomorphism.
Removing the point mapped to zero proves the topological assertion.  The
semigroup law follows directly from
\[
 \wt\pi_+(a,z)\wt\pi_+(a',z')
 =\wt\pi_+(aa',az'+z).
\]

Finally, the source proves
$\Phi_{\mathcal A(\pi_+)}\cong\wt G$ for the uniform closure
$\mathcal A(\pi_+)$.  Thus equality of the two spectra is exactly spectral
naturality of $\pi_+$.
\end{proof}

\section{Verification, limitations, and novelty}

\paragraph{Internal verification.}
The proof was checked against the source's Walter equation, its polar
decomposition statement, and the explicit formula for $\pi_+$.  The two
places where hidden regularity could enter were isolated.  The positive
classification uses only interval projections, positivity, and strong limits;
it never assumes an integral kernel for $T$.  The unitary classification uses
only bounded imaginary powers of the ratio operator; no subtraction of
unbounded logarithms is performed.  The resulting parameter formula agrees
with the polar formula displayed by Spronk--Stokke:
$z=b+iac$ means $c=a^{-1}\operatorname{Im}z$.

\paragraph{Scope limitation.}
This settles the explicit conjecture in the $ax+b$ example and hence one
important case of the paper's Conjecture~3.9(ii), because
$\pi_+\leq_q\lambda_G$.  It does not prove that every representation
quasi-contained in a left regular representation is spectrally natural.

\paragraph{Bounded novelty search.}
On 26 August 2026, the four run indexes and the local arXiv source corpus were
searched for arXiv:1112.4878, ``spectrally natural'', ``Wiener--Pitt
representation'', and the $\pi_+$/$ax+b$ character-space conjecture.  Web and
arXiv searches used the exact title, exact conjecture phrase, the terms
``spectrally natural'', ``$ax+b$ coefficient algebra spectrum'', and close
variants.  They found the source article, its journal record, citations that
use its general compactification framework, and no later paper stating a
resolution of this conjecture.  This is a bounded rather than exhaustive
novelty check.

\paragraph{Human review recommendation.}
Send to a harmonic analyst or operator algebraist.  The highest-value checks
are Lemma~\ref{lem:positive}'s passage from separated interval blocks to the
multiplication masa, and Lemma~\ref{lem:unitary}'s imaginary-power functional
calculus.  Neither step currently has an identified gap.

\begin{thebibliography}{9}

\bibitem{SpronkStokke}
N.~Spronk and R.~Stokke,
\emph{Matrix coefficients of unitary representations and associated compactifications},
Indiana Univ. Math. J. \textbf{62} (2013), no.~1, 99--148;
arXiv:1112.4878.

\end{thebibliography}

\end{document}
